\documentclass{article}
\pdfpagewidth=8.5in
\pdfpageheight=11in
\usepackage{ijcai19}

% Use the postscript times font!
\usepackage{times}
\usepackage{soul}
\usepackage{url}
\usepackage[hidelinks]{hyperref}
\usepackage[utf8]{inputenc}
\usepackage[small]{caption}
\usepackage{graphicx}
\usepackage{amsmath}
\usepackage{booktabs}
\usepackage{multirow}

\urlstyle{same}

\title{Deep learning project - RAG for Biomedical Question Answering}

\author{
Milka Gruber (27242045)
}

\begin{document}

\maketitle

\section{Introduction}
The objective of this project was to implement and evaluate a retrieval-augmented
generation (RAG) pipeline for biomedical multiple-choice question answering. The
central question was whether external information retrieved from a medical textbook corpus can
improve the answers of a relatively small open-source large language model (LLM).

The experiment uses \path{Qwen/Qwen2.5-1.5B-Instruct} as the LLM baseline,
\path{sentence-transformers/all-MiniLM-L6-v2} as the dense retriever, the
\path{MedRAG/textbooks} dataset as the external knowledge base, and the validation split
of \path{openlifescienceai/medmcqa} as the evaluation set. For this project
the dataset MedMCQA was chosen, as it is a dataset of multiple-choice questions,
which make the correctness of the results easier to evaluate. 


\section{RAG pipeline}

\subsection{Data and experiment configuration}

All the parameters used in the RAG pipeline are first read from a JSON configuration file.
They include the embedding and generation models, the number of questions and
textbook chunks used, the tested $k$ values, retrieval-query format, input-token limit,
random seed, and output paths. 

The corpus is formed from the training split of \texttt{MedRAG/textbooks}.
The full dataset consists of articles from Wikipedia, PubMed, StatPears and textbooks.
For this project I choose to use textbooks only, as it is the smallest part of the dataset.

The training split of the dataset is already preprocessed into chunks. 
Each chunk consists of an identifier, title, content and a text combined of title and the content. 
There are 125847 chunks in this dataset. 

Questions for evaluation come from the MedMCQA validation split of size 4183 questions.
Each entry consists of an identifier, question, four possible answers, 
the correct numerical answer to the question, choice type, explanation, subject name and the name of the topic.
All questions have a single answer reported as the correct one, so we treat the entire set
as a single answer type of question.

\subsection{Embeddings}

The text of every textbook chunk is encoded by all-MiniLM-L6-v2,
which is a sentence-transformer model that embedds short paragraphs
into a 384-dimensional vector space. 
The basic RAG pipeline consists of embedding both the textbook chunks
and the question query. Top k most similar documents to the question query
are retrieved and passed to the LMM, which generates the final answer.

Both textbook chunks and question query embeddings are normalized, 
so that similarity can be computed. These document embeddings are stored in
a FAISS \texttt{IndexFlatIP} index, which is a data structure that stores high-dimensional
vectors and optimizes for fast similarity search and nearest-neighbour retrieval.

Building embeddings for the complete corpus is expensive, so the FAISS index is
saved to disk (alongside it's metadata) and reused. 
The saved index is loaded only when the metadata matches the current configuration,
otherwise it is rebuilt.

The inner product between normalized vectors is equivalent to cosine similarity.
Thus, for a query vector $q$ and document vector
$d_i$, retrieval ranks documents using
\begin{equation}
    s(q,d_i)=q^\top d_i=\cos(q,d_i).
\end{equation}

\subsection{Retrieval}
I experimented with two retrieval query variants which are compared to the
embedded documents for the most similar document retrieval.
The first uses only the question text.
The second concatenates the question and the text of all four answer options, but
does not include the option labels A--D. For efficiency, the largest requested
$k$ is retrieved only once per question and smaller-$k$ experiments use prefixes of
that ranked list. 


\subsection{Prompting and generation}

The baseline prompt contains the question, four explicitly labelled options, and
the instruction to answer with only one letter. The RAG prompt prepends the top-$k$
retrieved chunks under a \emph{Context} heading and tells the model to use this
textbook context. Both variants use the same system message: \emph{You are a
medical question answering assistant.}

Qwen is placed in evaluation mode so that inference runs without gradient calculation.
The number of newly generated tokens is limited to 5.
For these experiments, the configured maximum input length was 32,768 tokens.
When the answer is generated, the original prompt tokens are removed from the
sequence, and a regular expression extracts the first standalone letter (A, B, C or D).

\subsection{Evaluation}

Accuracy is computed for the baseline and independently for every value of $k$.
This means that for each question and every k, separate answers are generated 
based on the different number of documents that are passed to the LLM. 

To show how retrieval changes model behaviour, the program additionally counts
\emph{fixed} answers (baseline incorrect, RAG correct) and \emph{broken} answers
(baseline correct, RAG incorrect). This distinction is useful because equal
accuracies do not report enough information on changes in behaviour. 

Every run is appended to JSON and plain-text result files.
The JSON output preserves information per question:
the correct answer, baseline and RAG predictions, outcome
category, retrieval query, ranked passages, titles, similarity scores, and dataset
metadata.

\section{Results}
We first ran the evaluation experiment on the entire 4183 MedMCQA validation set in 
two retrieval modes (queries with and without answers to the question). 
The results are shown in the Table~\ref{tab:full-results}. 
The closed-book Qwen baseline answered 1,833 questions correctly.
Overall the query with question and provided possible answers produced
overall better results. 
The best tested RAG setting for the question and answer query was $k=15$,
which answered 1,967 correctly and improved absolute accuracy by 3.20 percentage points.
The trends in both cases show that one retrieved document was insufficient
for the model to generate the correct answer, in which case it actually hurt
the performance. The performance in both query modes generally improved 
as more context was supplied. 


\begin{table}[!htb]
\centering
\small
\setlength{\tabcolsep}{3pt}
\resizebox{\linewidth}{!}{%
\begin{tabular}{lrrrrrrrr}
\toprule
& \multicolumn{4}{c}{Question + answer options} &
\multicolumn{4}{c}{Question only} \\
\cmidrule(lr){2-5}\cmidrule(lr){6-9}
Method & Correct & Accuracy & Fixed & Broken &
Correct & Accuracy & Fixed & Broken \\
\midrule
Baseline     & 1833 & 43.82\% & -- & -- & 1833 & 43.82\% & -- & -- \\
RAG ($k=1$)  & 1800 & 43.03\% & 522 & 555 & 1789 & 42.77\% & 447 & 491 \\
RAG ($k=3$)  & 1902 & 45.47\% & 608 & 539 & 1864 & 44.56\% & 522 & 491 \\
RAG ($k=5$)  & 1939 & 46.35\% & 630 & 524 & 1877 & 44.87\% & 556 & 512 \\
RAG ($k=8$)  & 1958 & 46.81\% & 635 & 510 & 1912 & 45.71\% & 563 & 484 \\
RAG ($k=10$) & 1943 & 46.45\% & 632 & 522 & 1907 & 45.59\% & 564 & 490 \\
%RAG ($k=12$) & 1964 & 46.95\% & 662 & 531 & X & X\% & X & X \\
RAG ($k=15$) & 1967 & 47.02\% & 660 & 526 & 1906 & 45.57\% & 590 & 517 \\
%RAG ($k=18$) & X & X\% & X & X & 1918 & 45.85\% & 588 & 503 \\
%RAG ($k=20$) & X & X\% & X & X & 1925 & 46.02\% & 583 & 491 \\
\bottomrule
\end{tabular}%
}
\caption{Results on the complete MedMCQA validation split for both retrieval-query
modes. ``Fixed'' and ``Broken'' are measured relative to the closed-book baseline.}
\label{tab:full-results}
\end{table}


\section{Retrieval analysis}
The results support the main hypothesis: relevant external context can improve a
small instruction-tuned model on biomedical questions, however this is too early to claim general improvement.
We see that quite a large number of answers which were initially correctly 
answered by Qwen, were later corrupted by RAG, which indicates that retrieval sometimes hurts correctness. 
In order to improve the pipeline we have to diagnose the reason the answer got broken.

I re-ran the entire pipeline for both retrieval cases on a smaller (n=500) random subsample of questions 
and saved the retrieved passages for $k=15$. 
I analyzed the retrieved passages manually and categorized the reason for failure in the following categories:
\begin{itemize}
    \item irrelevant retrieval: the retrieved passages were not related to the information needed for the question,
    \item relevant but insufficient retrieval: the passages were related to the topic but did not contain enough information to answer correctly,
    \item model ignored useful context: the correct information was retrieved, but the model did not use it,
    \item context supports wrong option: the retrieved passages contained information that made an incorrect option seem correct,
    \item ambiguous question / questionable correct answer in the dataset.
\end{itemize}

The reasons for fixed answers were categorized as follows:
\begin{itemize}
    \item direct answer found: the retrieved passages contained the correct answer,
    \item useful supporting evidence: the retrieved passages contained information that helped the model to choose the correct answer,
    \item multiple passages combined: the model combined information from multiple passages to arrive at the correct answer,
    \item context helped reject a wrong option: the retrieved passages contained information that helped the model to reject an incorrect option,
    \item correct despite weak retrieval: the model answered correctly even though the retrieved passages were not very useful,
    \item ambiguous question / questionable correct answer in the dataset.
\end{itemize}

The run where query contained both the question and answers for k=15 on 500 evaluation question set produced 
a 41\% accuracy for Qwen and 47.20\% accuracy for Qwen with RAG pipeline. 
The RAG pipeline fixed 89 previously incorrect answers and broke 58 previously correct answers. 

The run with question only query for k=15 on the same 500 evaluation question set as in the first case produced 
a 41\% accuracy for Qwen and 44.20\% accuracy for Qwen with RAG pipeline. 
The RAG pipeline fixed 72 previously incorrect answers and broke 56 previously correct answers. 

\subsection{Analysis of corrupted answers}

For each corrupted answer, all 15 retrieved passages were examined and one primary
failure category was assigned (the categories are mutually exclusive).
Table~\ref{tab:corruption-analysis} summarizes the analysis.

\begin{table}[!htb]
\centering
\small
\setlength{\tabcolsep}{3pt}
\resizebox{\linewidth}{!}{%
\begin{tabular}{lrrrr}
\toprule
& \multicolumn{2}{c}{Question + answers} &
\multicolumn{2}{c}{Question only} \\
\cmidrule(lr){2-3}\cmidrule(lr){4-5}
Primary failure category & Count & Share & Count & Share \\
\midrule
Irrelevant retrieval                         & 11 & 19.0\% & 16 & 28.6\% \\
Relevant but insufficient retrieval          & 17 & 29.3\% & 16 & 28.6\% \\
Model ignored useful context                  & 14 & 24.1\% & 16 & 28.6\% \\
Context supports wrong option                 & 11 & 19.0\% &  4 &  7.1\% \\
Ambiguous question / questionable gold label  &  5 &  8.6\% &  4 &  7.1\% \\
\midrule
Total                                         & 58 & 100\%  & 56 & 100\%  \\
\bottomrule
\end{tabular}%
}
\caption{Primary causes assigned to corrupted answers in the two 500-question
experiments at $k=15$.}
\label{tab:corruption-analysis}
\end{table}

The corruption analysis shows that the two query modes produced a similar number
of broken answers: 58 for the question and answer query and 56 for the
question-only query. In both modes, approximately 69\% of corrupted answers had
no retrieved passage containing the information required for the correct answer
(40 of 58 and 39 of 56, respectively).

Further inspection suggested a corpus-coverage problem for dental questions.
Dental questions accounted for 20 of the 58 corrupted answers. In 16 of these
20 cases (80\%), the retrieved passages were either irrelevant or related to the
topic but insufficient to answer the question. The retriever often returned
general medical information or broad passages about teeth that did not contain
the specialized dental facts required to distinguish between the answer options.
%We identify the lack of information
%that would inform the correct answer as the main reason for failure in both query modes.

The distribution of types of failures differed between the two query modes. With question-only
retrieval, the three most common problems were irrelevant passages, insufficient
information, and the model ignoring useful passages. Each occurred in 16 of 56
cases. Adding the answer options reduced irrelevant retrieval from
28.6\% to 19.0\% and ignored useful context from 28.6\% to 24.1\%. However, it
also increased cases where the passages supported a wrong answer from 7.1\% to
19.0\%. Therefore, answer options can help retrieve more relevant information,
but they can also direct the search toward an incorrect option. Overall, using
the question and answer options produced higher accuracy, despite this risk.

\subsection{Analysis of fixed answers}

\begin{table}[!htb]
\centering
\small
\setlength{\tabcolsep}{3pt}
\resizebox{\linewidth}{!}{%
\begin{tabular}{lrrrr}
\toprule
& \multicolumn{2}{c}{Question + answers} &
\multicolumn{2}{c}{Question only} \\
\cmidrule(lr){2-3}\cmidrule(lr){4-5}
Primary reason for correction & Count & Share & Count & Share \\
\midrule
Direct answer found                              & 37 & 41.6\% & 23 & 31.9\% \\
Useful supporting evidence                       &  9 & 10.1\% &  6 &  8.3\% \\
Multiple passages combined                       &  2 &  2.2\% &  1 &  1.4\% \\
Context helped reject a wrong option              &  9 & 10.1\% & 10 & 13.9\% \\
Correct despite weak retrieval                    & 30 & 33.7\% & 30 & 41.7\% \\
Ambiguous question / questionable gold label      &  2 &  2.2\% &  2 &  2.8\% \\
\midrule
Total                                             & 89 & 100\% & 72 & 100\%  \\
\bottomrule
\end{tabular}%
}
\caption{Primary reasons assigned to fixed answers in the two 500-question
experiments at $k=15$.}
\label{tab:fixed-analysis}
\end{table}

The analysis of fixed answers also shows a difference between the two query
modes. The question-and-answer query corrected 89 answers, while the
question-only query corrected 72. With answer options included, direct answers
were found in the retrieved passages in 37 cases (41.6\%), compared with 23 cases
(31.9\%) for question-only retrieval. This suggests that answer options helped
the retriever find passages containing the required information.

However, 30 answers were corrected despite weak retrieval in both modes. This
represented 33.7\% of corrections for the question-and-answer query and 41.7\%
for the question-only query. In these cases, the model probably relied more on
its existing knowledge or guessed correctly than on the retrieved passages.
The remaining corrections came from supporting evidence, combining passages,
or using context to reject a wrong option. Overall, including the answer options
produced more corrected answers and more cases where the correct answer was
directly supported by the retrieved text.

\section{Further work}
Based on the analysis results, we can determine that the query mode with
both the question and answers is more beneficial to the RAG pipeline, 
so any further work shall focus on this query mode. 

We identified the lack of information that would inform the correct answer
as the main reason for failure in both query modes. This could be fixed by improving the corpus coverage,
for example by adding wikipedia articles from the same dataset.

The issue where the retrieved passages confused the model into choosing a wrong answer
could be helped by reranking the retrieved passages based on their relevance to the question and answer options.

Trying out different LMMs and embedding models could also improve the results. 
Because specific words in medicine can have different meanings across subfields, it is 
important for the embedding model to be able to determine the context of the words to understand the correct meaning. 
This could be achieved by using a domain-specific embedding model.




\end{document}
